\subsection{QUOTIENT SPACES}

DEFINITION 3. Let X be a topological space, R an equivalence relation on X. The quotient space of X by R is the quotient set X/R with the topology which is the quotient of the topology of X by the relation R (§ 2, no. 4, Example 1).

Unless the contrary is expressly stated, whenever we speak of X/R as a topological space, it is to be understood that we mean the quotient space of X by R. We shall often say that this topological space is the space obtained by identifying the points of X which belong to the same equivalence class mod R.

Let \(\varphi\) be the canonical mapping \(X \to X/R\) . By definition (§ 2, no. 4, Proposition 6 and its corollary) the open (resp. closed) sets in \(X/R\) are the sets A such that \(\bar{\varphi}^{1}(A)\) is open (resp. closed) in \(X\) ; in other words, the open (resp. closed) sets in \(X/R\) are in one-to-one correspondence with the open (resp. closed) subsets of \(X\) which are saturated with respect to \(R\) and are the canonical images of these subsets.

PROPOSITION 6. Let X be a topological space, R an equivalence relation on X, φ the canonical mapping of X onto X/R; then a mapping f of X/R into a topological space Y is continuous if and only if f \(\circ\) φ is continuous on X.

This is a particular case of § 2, no. 4, Proposition 6; it expresses the fact that the quotient topology is the final topology for the mapping \(\varphi\) .

Proposition 6 shows that there is a one-to-one correspondence between the continuous mappings of X/R into Y and the continuous maps of X into Y which are constant on each equivalence class mod R.

Example. * Consider the equivalence relation \(x \equiv y \pmod{1}\) on the real line \(\mathbf{R}\) ; the quotient space of \(\mathbf{R}\) by this relation is called the one-dimensional torus and is denoted by \(\mathbf{T}\) . The equivalence class of a point \(x \in \mathbf{R}\) consists of all the points \(x + n\) , where \(n\) runs through the set \(\mathbf{Z}\) of rational integers. By Proposition 6 there is a one-to-one correspondence between the continuous functions on \(\mathbf{T}\) and the continuous functions on \(\mathbf{R}\) which are periodic with period 1. We shall return to this important example in Chapter V, § 1. *

COROLLARY. Let X, Y be two topological spaces, R (resp. S) an equivalence relation on X (resp. Y), and let \(f: X \to Y\) be a continuous mapping which is compatible with the equivalence relations R and S (Set Theory R, § 5, no. 8); then the mapping \(g: X/R \to Y/S\) induced by \(f\) (Set Theory R, § 5, no. 8) is continuous.

This is a particular case of a general property of quotient structures (Set Theory, Chapter IV, § 2, no. 6, criterion CST 20).

PROPOSITION 7 (Transitivity of quotient spaces). Let R and S be two equivalence relations on a topological space X such that R implies S, and let S/R be the quotient equivalence relation on the quotient space X/R (Set Theory R, § 5, no. 9). Then the canonical bijection \((\mathbf{X}/\mathbf{R})/(\mathbf{S}/\mathbf{R}) \to \mathbf{X}/\mathbf{S}\) is a homeomorphism.

This is a particular case of the transitivity of final topologies (§ 2, no. 4, Proposition 7. Cf. Set Theory, Chapter IV, § 2, no. 3, criterion CST 21).

